\section{Two double covers and their common fourfold cover}\label{sec:square}

Let $C$ be a closed smooth complex curve of genus $g\ge2$, and let $N$
have complex rank two and degree $2g-2$. For a real Cauchy--Riemann
operator $D$ with Dolbeault symbol on $N$, let $D_u$ be its pullback
to $u:X\to C$. On a nodal domain the sections on the normalization
have matching values. We call $D$ \emph{rigid through degree $M$} if
$D_u$ is injective for every connected stable map of degree between
one and $M$. At such a parameter write
\begin{equation}\label{eq:fixed}
L^C_{d,k}(D)=\int_{[\overline{\mathcal M}_{k}(C,d)]^{\rm vir}}
            e(\operatorname{coker}D_u),\qquad
 L^C_d(D;v)=\sum_k L^C_{d,k}(D)v^{2k-2}.
\end{equation}
We orient the normal complex by continuation from complex-linear
operators. Its real cokernel has rank $4(k-1)-4d(g-1)$, the virtual
real dimension of the map space. The fixed-target comparison theorem
identifies the integral with the contribution of the isolated
zero section in the homogeneous normal bundle
\cite[Theorem~1.2]{zinger}. For its balanced product we write $P_C(D;q)$;
its construction and degreewise uniqueness are given below.
Set $H_d=1+d(g-1)$. The sign of an invertible operator on an unramified
cover $T$ will be written $\epsilon_T(D)$.

Choose distinct nonzero classes $\alpha_1,\alpha_2\in H^1(C;\mathbb F_2)$.
Write $Y_i$ for the double classified by $\alpha_i$. Since the characters
are independent, $Z=Y_1\times_CY_2$ is connected. Each corner of a
two-parameter family will be denoted by the signs of its two parameters.
In the following displays a count at nonintegral degree is zero.
\begin{theorem}\label{thm:square}
There is a family on $(C,N)$ for which, at every finite cutoff $M$,
a sufficiently small rectangle consists of four rigid open quadrants
and the two crossing axes. For $d\le M$ and every $k$, the connected
difference satisfies
\[
 L_{++}-L_{-+}-L_{+-}+L_{--}
      = L^{Z}_{d/4,k}(D_{++}|_Z).
\]
Here $L_\epsilon=L^C_{d,k}(D_\epsilon)$. The quotient of the corresponding
balanced products is
\[
 \frac{P_{++}P_{--}}{P_{-+}P_{+-}}=P_Z(D_{++}|_Z;q^4)
          \pmod {q^{M+1}} .
\]
For $M\ge4$, the first differing exponent has label four, elementary
genus $4g-3$, and value $\epsilon_Z(D_{++})$.
\end{theorem}

The geometric reason for the factor one is worth spelling out. Embed
the sum $W$ of the two flat sign lines in $N_{\mathbb R}$. Indeed,
$W\oplus W$ is oriented with $w_2=w_1(W)^2=0$; the oriented rank-four
bundle $N_{\mathbb R}$ also has $w_2=0$ and is isomorphic to it on
a surface. Prescribe the parallel vectors in these copies to be
central kernel sections. Pair a variation of the normal operator with
the first twisted cokernel while killing the second prescribed
section, and do the converse for the other parameter. The construction
in \S\ref{sec:family} makes these the only kernels apart from their
pullbacks and restrictions. At a stable map each sign reverses exactly
when its character pulls back trivially, including around loops through
nodes. On taking the alternating sum, a map sees a multiplier four
if it lifts to both doubles and zero otherwise. A specified map in
the surviving locus has four lifts to $Z$ with its domain fixed.
Representable finite \'etale pushforward divides its upstairs integral
by four, giving the formula. The orientation and compactness statements
needed for this computation are proved in
\S\S\ref{sec:family}--\ref{sec:oriented}.

More explicitly the increments from left to right on the two opposite
edges are the positive-side counts in degree $d/2$ on $Y_1$, with the
second sign positive or negative respectively. Subtracting them is
the sign-wall calculation for $Z\to Y_1$. Subtracting the other pair
instead uses $Z\to Y_2$ and gives the identical answer. The oriented
sum along the entire boundary is zero. At degree four and genus
$4g-3$ the only contributing common lift downstairs is $Z$ itself,
of deck order four. Its four alternating masses give its sign.
In the product this also follows from the signed degree-one contribution
on $Z$, which is the full series
$\epsilon_Z(D_{++})(2\sin(v/2))^{2(4g-3)-2}$.
